% ============================================================
%  SCHOLA DOMESTICA — Level 1, Week 3, Lesson 12
%  Topic: Counting Backward — From 5 to 0
%  8–10 problems | No speed drill
%  Compile with XeLaTeX
% ============================================================
\documentclass[12pt,letterpaper]{article}
\usepackage{sd_style}

\renewcommand{\lessonnum}{12}
\renewcommand{\lessonweek}{3}

\begin{document}

\lessonheader{3}{12}{Counting Backward from 5 to 0}

\begin{instructbox}
\noindent\textbf{\color{sd-blue}Instruction} \textit{(read aloud):}\\[4pt]
We know how to count \textit{forward}: 0, 1, 2, 3, 4, 5.
Today we learn to count \textit{backward}: 5, 4, 3, 2, 1, 0.
Counting backward means going one \textit{less} each time.
It is like a rocket launch — 5, 4, 3, 2, 1, 0 \ldots\ blast off!
\end{instructbox}

\vspace{8pt}
% Backward number line
\noindent\textbf{\color{sd-blue}Count backward on the number line:}
\vspace{8pt}

\begin{center}
\begin{tikzpicture}
  \draw[sd-blue, very thick, <-] (-0.3,0) -- (5.8,0);
  \foreach \n in {0,...,5}{
    \draw[sd-blue, thick] (\n, -0.15) -- (\n, 0.15);
    \node[below] at (\n, -0.25) {\textbf{\n}};
  }
  % Arrow showing direction
  \draw[sd-gold, very thick, ->] (5.2, 0.5) -- (4.2, 0.5);
  \node[above] at (4.7, 0.55) {\color{sd-gold}\small\textit{this way}};
\end{tikzpicture}
\end{center}

\goldrule
\noindent\textbf{\color{sd-blue}Practice}
\vspace{6pt}

\prob Fill in the blanks — count backward.\\[6pt]
\hspace{1cm}{\Large 5 \quad \ansblank \quad 3 \quad \ansblank \quad 1 \quad \ansblank}

\vspace{10pt}
\prob What comes just \textbf{before} 4 when counting backward?\\[4pt]
\hspace{1cm}5 \quad \ansblank \quad 3

\vspace{10pt}
\prob What comes just \textbf{before} 2 when counting backward?\\[4pt]
\hspace{1cm}3 \quad \ansblank \quad 1

\vspace{10pt}
\prob \textit{(Teacher reads)} Five frogs sat on a log.
One jumped into the pond. Count backward to find how many are left.\\[4pt]
\hspace{1cm}Start at: 5 \quad Count back 1 \quad Answer: \ansblank

\vspace{10pt}
\prob Count backward and write the missing numbers:\\[6pt]
\begin{center}
\begin{tikzpicture}
  \foreach \n/\v in {1/5, 2/{\ }, 3/3, 4/{\ }, 5/1, 6/0}{
    \draw[sd-blue, thick] (\n*1.5-1.5, 0) rectangle (\n*1.5-0.3, 0.9);
    \ifnum\n=2
      \node at (\n*1.5-0.9, 0.45) {\Large\ansblank};
    \else\ifnum\n=4
      \node at (\n*1.5-0.9, 0.45) {\Large\ansblank};
    \else
      \node at (\n*1.5-0.9, 0.45) {\Large\textbf{\v}};
    \fi\fi
  }
  \draw[sd-gold, very thick, ->] (6.6, 0.45) -- (7.2, 0.45);
  \node[right] at (7.2, 0.45) {\color{sd-gold}\small\textit{blast off!}};
\end{tikzpicture}
\end{center}

\vspace{10pt}
\prob Write the number that is \textbf{one less} than each.\\[6pt]
\begin{tabular}{ccccc}
one less than 5: & one less than 3: & one less than 1: & one less than 4: & one less than 2:\\[4pt]
\ansblank & \ansblank & \ansblank & \ansblank & \ansblank
\end{tabular}

\vspace{10pt}
\prob \textit{(Teacher reads)} A candle burned down from 4 inches tall.
Each day it got 1 inch shorter.
Count backward from 4 to show the heights: \quad
4 \quad \ansblank \quad \ansblank \quad \ansblank \quad \ansblank

\vspace{10pt}
\prob Count backward aloud from 5 to 0 and mark each number as you say it.\\[6pt]
\begin{center}
\begin{tabular}{cccccc}
\Large\textbf{5} & \Large\textbf{4} & \Large\textbf{3} & \Large\textbf{2} & \Large\textbf{1} & \Large\textbf{0}\\[4pt]
$\square$ & $\square$ & $\square$ & $\square$ & $\square$ & $\square$
\end{tabular}
\end{center}

\goldrule

\begin{checkbox}
\noindent\textbf{\color{sd-green}Knowledge Check}\\[4pt]
\setcounter{probnum}{0}
\prob Count backward from 5: \quad 5, \ansblank, \ansblank, \ansblank, \ansblank, \ansblank\\[8pt]
\prob What is one less than 3? \quad\ansblank
\end{checkbox}

\vfill
\noindent{\color{sd-blue}$\bigstar$ \textbf{Enrichment — Tally Marks:}}\\
People have counted with tally marks for thousands of years.
Every four marks get crossed by a fifth: \quad
\begin{tikzpicture}[baseline=0.1cm]
  \foreach \x in {0,...,3}\draw[sd-blue,very thick](\x*0.3, 0)--(\x*0.3, 0.8);
  \draw[sd-blue,very thick,rotate around={-70:(0.45,0.4)}](0,0.4)--(0.9,0.4);
\end{tikzpicture}
\quad = 5\\[4pt]
Count backward using tally marks: draw 5, cross one off at a time.

\vspace{4pt}\hrule\vspace{3pt}
{\footnotesize\color{gray}\textit{Teacher note: The rocket launch countdown (5–4–3–2–1–0) is an excellent mnemonic.
  Practice orally several times. If student confuses "one less" direction,
  return to the physical number line and hop left one space at a time.}}

\end{document}
